\chapter{Zagro\.{z}enia i~techniki manipulacji multimediami}
\label{ch:zagrozenia}

Rozdzia\l{} porz\k{a}dkuje poj\k{e}cia i~mechanizmy potrzebne do interpretacji
cz\k{e}\'{s}ci implementacyjnej. W~odróżnieniu od~rozdziału~\ref{ch:tlo},
który syntetyzuje literatur\k{e}, nacisk po\l{}o\.{z}ono tu na~operacyjny kontekst
znakowania wodnego i~proweniencji tre\'{s}ci.

% ===================================================================
\section{Taksonomia technik manipulacji}
\label{sec:taksonomia}

Techniki generowania i~manipulacji multimediami mo\.{z}na podzieli\'{c}
na~pi\k{e}\'{c} g\l{}\'{o}wnych kategorii:

\begin{enumerate}
  \item \textbf{Podmiana twarzy (face swap)} ---
        zamiana twarzy jednej osoby na~twarz innej przy
        zachowaniu ruchu g\l{}owy i~mimiki orygina\l{}u.
        Historycznie realizowana przez sieci GAN
        (DeepFaceLab, FaceSwap); wsp\'{o}\l{}cze\'{s}nie
        tak\.{z}e przez modele dyfuzyjne.

  \item \textbf{Synteza twarzy (face reenactment)} ---
        przeniesienie ekspresji i~ruchu ust z~nagrania
        \'{z}r\'{o}d\l{}owego na~docelowe nagranie innej osoby.
        Techniki: First Order Motion Model, DiffusionFace.

  \item \textbf{Synteza ca\l{}ej sceny (text-to-video)} ---
        generowanie pe\l{}nej sekwencji wideo z~opisu
        tekstowego (OpenAI Sora, Runway Gen-3/4, Kling AI,
        Luma Dream Machine, Pika Labs).
        U\.{z}ywa architektur Diffusion Transformer (DiT)
        lub dyfuzji latentnej.

  \item \textbf{Manipulacja audio (voice cloning)} ---
        klonowanie g\l{}osu na~podstawie pr\'{o}bki kilku sekund
        nagrania (ElevenLabs, Tortoise TTS).
        Cz\k{e}sto \l{}\k{a}czone z~manipulacj\k{a} wideo
        w~kompletnych \textit{deepfake}\'{a}ch audio-wideo.

  \item \textbf{Manipulacja dokumentami i~zdj\k{e}ciami} ---
        edycja tre\'{s}ci obrazu (inpainting) lub dokumentu
        (paszporty, certyfikaty) przy u\.{z}yciu modeli
        generatywnych.
\end{enumerate}

Skal\k{e} zagro\.{z}enia ilustruje fakt, \.{z}e raport Europolu z~2023~r.
wskazuje na bardzo szybki wzrost udziału treści generowanych lub modyfikowanych
przez AI w~Internecie~\cite{Europol2023Deepfakes}.
Wp\l{}yw dezinformacji multimedialnej dokumentuj\k{a} liczne
badania, w~szczeg\'{o}lno\'{s}ci w~kontekstach wyborczym,
bankowym i~wojennym~\cite{Khan2025ComprehensiveSurvey}.

% ===================================================================
\section{Znakowanie wodne tre\'{s}ci generowanych przez AI}
\label{sec:watermarking_detail}

Znakowanie wodne (\textit{digital watermarking}) polega na~osadzeniu
w~pliku multimedialnym informacji identyfikacyjnej w~taki spos\'{o}b,
aby by\l{}a ona trudna do~usuni\k{e}cia lub niezauwa\.{z}alna
dla obserwatora~\cite{Luo2025WatermarkingAIGC}.
W~kontek\'{s}cie tre\'{s}ci generowanych przez AI technologia ta
pozwala na~jednoznaczne oznaczenie materia\l{}u ju\.{z} w~momencie
jego tworzenia, co~tworzy alternatywn\k{a} \'{s}cie\.{z}k\k{e}
weryfikacji, niezale\.{z}n\k{a} od~analizy artefakt\'{o}w wizualnych.

% -------------------------------------------------------------------
\subsection{Znakowanie wodne widoczne}
\label{sec:watermark_visible}

Widoczny znak wodny (\textit{visible watermark}) jest celowo
prezentowany w~polu widzenia odbiorcy --- typowo jako logotyp,
napis lub wzorzec geometryczny na\l{}o\.{z}ony na~tre\'{s}\'{c}.
W~praktyce pe\l{}ni on g\l{}\'{o}wnie funkcj\k{e} informacyjn\k{a}
lub licencyjn\k{a}, a~nie zabezpieczaj\k{a}c\k{a}. Jego zalet\k{a}
jest prosta weryfikacja przez OCR, dopasowanie szablonowe albo analiz\k{e}
typowych obszar\'{o}w ekranu, natomiast wad\k{a} pozostaje niska odporno\'{s}\'{c}
na przycinanie, kadrowanie i~inpainting.

% -------------------------------------------------------------------
\subsection{Znakowanie wodne niewidoczne}
\label{sec:watermark_invisible}

Niewidoczny znak wodny (\textit{invisible\,/\,steganographic watermark})
zostaje osadzony w~danych pliku w~spos\'{o}b niezauwa\.{z}alny
go\l{}ym okiem, ale potencjalnie odczytywalny przez specjalistyczny
dekoder.

\subsubsection{Metody osadzania}
\label{sec:watermark_embed}

Najcz\k{e}\'{s}ciej spotykane klasy rozwi\k{a}za\'{n} to:
\textbf{LSB}, techniki w~dziedzinie cz\k{e}stotliwo\'{s}ci
(\textbf{DCT/DWT}) oraz metody neuronowe, takie jak
\textbf{HiDDeN} i~\textbf{StegaStamp}
~\cite{Zhu2018HiDDeN,Fernandez2024BlindWatermark}.
Dla wideo rozwa\.{z}a si\k{e} tak\.{z}e osadzanie sygna\l{}u
w~wymiarze temporalnym, np.\ przez modyfikacj\k{e} fazy ruchu
lub percepcyjnie niewidoczne wzorce migotania
~\cite{Luo2025WatermarkingAIGC}.

\subsubsection{Metody wykrywania znaku niewidocznego}
\label{sec:watermark_detect}

W praktyce wyr\'{o}\.{z}nia si\k{e} detekcj\k{e} \textit{blind}
(z~kluczem, bez dost\k{e}pu do~orygina\l{}u), detekcj\k{e} referencyjn\k{a}
(\textit{non-blind}) oraz podej\'{s}cia po\'{s}rednie, oparte na analizie
residuum, szumu albo wzorc\'{o}w pokroju PRNU
~\cite{Wang2024ReliabilityPerspective}. Z~punktu widzenia niniejszej pracy
istotne jest to, \.{z}e pe\l{}ne odczytanie takiego znacznika wymaga zwykle
dedykowanego dekodera i~klucza, dlatego sama analiza tre\'{s}ci mo\.{z}e
co najwy\.{z}ej dostarcza\'{c} sygna\l{}\'{o}w proxy.

\subsubsection{Odporno\'{s}\'{c} a~ataki}
\label{sec:watermark_attacks}

Literatura wskazuje, \.{z}e niewidoczne znaki wodne pozostaj\k{a}
podatne na~kompresj\k{e}, ataki geometryczne, filtrowanie,
perturbacje adwersaryjne oraz regeneracj\k{e} przez modele generatywne
~\cite{Zhao2023InvisibleRemovable,Piet2024SoKWatermarking}. Nie nale\.{z}y
ich zatem traktowa\'{c} jako uniwersalnego i~samowystarczalnego mechanizmu
potwierdzania autentyczno\'{s}ci.

% -------------------------------------------------------------------
\subsection{Metryki jako\'{s}ci znakowania wodnego}
\label{sec:watermark_metrics}

Ewaluacja systemu znakowania wodnego wymaga oceny
dwóch przeciwstawnych w\l{}a\'{s}ciwo\'{s}ci.

\textbf{Niepostrzegalno\'{s}\'{c} (imperceptibility).}
Mierzona metrykami postrzegania jako\'{s}ci obrazu:
\textbf{PSNR} (\textit{Peak Signal-to-Noise Ratio}),
\textbf{SSIM} (\textit{Structural Similarity Index})
oraz \textbf{LPIPS} (\textit{Learned Perceptual Image Patch Similarity}).
Wy\.{z}szy PSNR i~SSIM oznacza mniejsz\k{a} ingerencj\k{e} w~obraz.

\textbf{Odporno\'{s}\'{c} (robustness).}
Mierzona \textbf{BER} (\textit{Bit Error Rate}) ---
odsetkiem b\l{}\k{e}dnie odczytanych bit\'{o}w wiadomo\'{s}ci po~ataku.
Cel: BER~$< 5\%$ po~typowych transformacjach
(kompresja, zmiana rozdzielczo\'{s}ci).

Dodatkowo stosuje si\k{e} metryki skuteczno\'{s}ci detekcji
w~uj\k{e}ciu binarnym:
\begin{itemize}
  \item \textbf{TPR} (\textit{True Positive Rate}) ---
        odsetek poprawnie wykrytych znak\'{o}w.
  \item \textbf{FPR} (\textit{False Positive Rate}) ---
        odsetek plik\'{o}w nieoznaczonych, kt\'{o}re system
        b\l{}\k{e}dnie sklasyfikowa\l{} jako oznakowane.
  \item \textbf{AUC-ROC} ---
        pole pod krzyw\k{a} klasyfikatora dla r\'{o}\.{z}nych
        prog\'{o}w detekcji.
\end{itemize}

% ===================================================================
\section{Content Provenance -- C2PA}
\label{sec:c2pa}

\subsection{Czym jest C2PA?}
\label{sec:c2pa_opis}

\textbf{C2PA} (\textit{Coalition for Content Provenance and Authenticity})
to~otwarta specyfikacja techniczna definiuj\k{a}ca
\textbf{Content Credentials}, czyli podpisane cyfrowo metadane opisuj\k{a}ce
histori\k{e} powstania i~edycji pliku multimedialnego~\cite{C2PA2024Spec}.
C2PA nie jest systemem detekcji artefakt\'{o}w, lecz mechanizmem
\textbf{proweniencji kryptograficznej}.

\subsection{Architektura i~dzia\l{}anie}
\label{sec:c2pa_architektura}

Plik zgodny z~C2PA mo\.{z}e zawiera\'{c} manifest z~informacjami
o~wykonanych operacjach, asercjach dotycz\k{a}cych narz\k{e}dzia lub modelu
oraz podpisie cyfrowym. Po stronie odbiorcy weryfikacja sprowadza si\k{e}
do~odczytu manifestu, sprawdzenia integralno\'{s}ci pliku, wa\.{z}no\'{s}ci
podpisu i~tre\'{s}ci asercji, np.\ \texttt{c2pa.ai.generative}.

\subsection{C2PA a~detekcja AI -- rola w~badaniu}
\label{sec:c2pa_detekcja}

W~niniejszym badaniu C2PA pe\l{}ni rol\k{e} sygna\l{}u wysokiej pewno\'{s}ci:
wa\.{z}ny manifest z~asercj\k{a} \texttt{c2pa.ai.generative} stanowi
mocn\k{a} przes\l{}ank\k{e} pochodzenia generatywnego. Ograniczeniem pozostaje
jednak niepowszechne wdro\.{z}enie standardu, mo\.{z}liwo\'{s}\'{c} utraty
metadanych podczas reenkodacji oraz konieczno\'{s}\'{c} walidacji \l{}a\'{n}cucha
zaufania. Brak manifestu nie mo\.{z}e by\'{c} wi\k{e}c interpretowany
jako dow\'{o}d autentyczno\'{s}ci.
